\chapter{\abstractname}

A language model's weights are large arrays of numbers, called tensors, that the model multiplies on every generated token. The full set of tensors lives on disk in a single file, and the model runs fastest when the whole file is loaded into RAM. RAM has become expensive. DDR5 contract prices have roughly doubled since early 2025, and AI accelerator demand is forecast to consume about 20\% of global DRAM supply in 2026. NVMe SSDs are cheap and fast. Our goal is to extend llama.cpp so that a large model can run on a machine that does not have enough RAM to hold the whole model.

We started our analysis by building a tensor-level tracker that records, for every operation the model performs, which tensor it reads and which region of memory the data comes from. The trace shows that only 26.91\% of GPT-OSS-20B's weights and 7.69\% of GPT-OSS-120B's weights are read per generated token, because each Mixture-of-Experts layer activates only 4 of its 32 experts on 20B, and 4 of 128 on 120B. The inactive experts, which make up most of the model file, can stay on the SSD. Our backend pins the always-needed weights (the attention, output, and norm tensors) in RAM so the operating system cannot evict them, and prefetches the expert weights from the SSD asynchronously, overlapping the disk reads with the matrix multiplications that consume earlier reads. GPT-OSS-20B (12.85~GiB on disk) runs at 14.32 tokens per second on our test machine when the whole model fits in RAM. With limited RAM, where the model no longer fits, the default llama.cpp loader drops to 1.44 tokens per second. Our backend, on the same setting, runs the same model at 5.36 tokens per second. GPT-OSS-120B (60.88~GiB on disk) does not fit in the 30~GiB of RAM on our test machine at all. With limited RAM the default loader runs it at 4.96 tokens per second and our backend at 8.06 tokens per second.
